%% kapa: exam=1 round=31 year=2020 period=2 form=A subject=appraisal_law subject_ko=감정평가관계법규 seq=7 number=7
%% source: Q-Net 감정평가사 기출문제 (cst00309) artlSeq=5207219 file=extracted/1008/5207219/A형/A형/2020년도 제31회 감정평가사 1차 시험_2교시_A형.pdf
%% conversion: kordoc 4.24.51 -> tools/gen1.py (auto); review=proofread-sonnet; flags=none
\begin{question}{7}
국토의 계획 및 이용에 관한 법령상 ‘법률 등의 위반자에 대한 처분’을 함에 있어서 청문을 실시해야 하는 경우로 명시된 것을 모두 고른 것은?
\begin{framedtext}
ㄱ. 개발행위허가의 취소 \par
ㄴ. 개발행위의 변경허가 \par
ㄷ. 토지거래계약 허가의 취소 \par
ㄹ. 실시계획인가의 취소 \par
ㅁ. 도시ㆍ군계획시설사업의 시행자 지정의 취소 \par
\end{framedtext}
\begin{choices}
  \item ㄱ, ㄴ
  \item ㄴ, ㄷ
  \item ㄱ, ㄹ, ㅁ
  \item ㄷ, ㄹ, ㅁ
  \item ㄱ, ㄴ, ㄷ, ㄹ
\end{choices}
\end{question}
\answer{3}
